\begin{figure}[htbp]
\centering
\begin{tikzpicture}[scale=1.05]
% 支割线：正实轴上一段，用红色粗线表示
\draw[very thick, red!70] (0.25,0) -- (2.7,0);
\node[red!70, above] at (1.7,0.02) {\small 支割线};

% 钥匙孔：上沿（R -> eps），从外面绕过支割线
\draw[very thick, -{Stealth}]
  (2.7,-0.14) .. controls (1.5,-0.14) and (0.5,-0.14) .. (0.25,-0.14);
% 左端小圆（逆时针，从下到上）
\draw[very thick] (0.25,0.14) arc (90:270:0.14);
% 下沿（eps -> R）
\draw[very thick, -{Stealth}]
  (0.25,0.14) .. controls (0.5,0.14) and (1.5,0.14) .. (2.7,0.14);
% 右端：从下沿末端经大圆回到上沿起点（逆时针沿大半圆）
\draw[very thick, -{Stealth}]
  (2.7,0.14) arc (4:340:2.7) ;

% 缺口标记
\draw (-3.0,0) -- (3.4,0) node[right] {$\mathrm{Re}\,z$};
\draw (0,-3.0) -- (0,3.2) node[above] {$\mathrm{Im}\,z$};

% 极点示例（在支割线外）
\fill[blue!70] (-1.2,1.2) circle (2pt) node[right] {\small $z=-1$};

\node at (2.95,-0.15) {\small $R$};
\node at (0.55,0.35) {\small $\varepsilon$};
\end{tikzpicture}
\caption{钥匙孔回路（keyhole contour）。沿支割线上方从 $R$ 走到 $\varepsilon$，绕小圆一圈，再沿下方从 $\varepsilon$ 走回 $R$，最后绕大半圆闭合。上下两侧的积分不抵消——支割线两侧函数取不同的分支，差一个 $e^{2\pi ia}$，这个差正是积分的来源。}
\end{figure}
